% Shared preamble for the in-class practice handouts and answer keys.
% Each main file sets \chapterlabel and \handoutkind; compile with xelatex.
\documentclass[UTF8, scheme=plain, fontset=fandol, 11pt, a4paper]{ctexart}

% Same CJK fonts as elegantbook-zh-tw.cls when installed (e.g. on Overleaf);
% otherwise fall back to the macOS Traditional Chinese fonts (Fandol lacks many TC glyphs).
\IfFontExistsTF{AR PL UMing TW MBE}{%
    \setCJKmainfont[ItalicFont=AR PL UKai TW, BoldFont=cwTeXHeiBold]{AR PL UMing TW MBE}%
}{\IfFontExistsTF{Songti TC}{%
    \setCJKmainfont[ItalicFont=Kaiti TC, BoldFont=Heiti TC]{Songti TC}%
}{}}

\usepackage[margin=2cm, top=2.2cm, bottom=2cm]{geometry}
\usepackage{amsmath, amssymb}
\usepackage{xcolor}
\usepackage{enumitem}
\usepackage{fancyhdr}
\usepackage[most]{tcolorbox}

\definecolor{structurecolor}{RGB}{60,113,183}
\definecolor{third}{RGB}{0,160,152}

\newcommand{\PP}{\mathbb{P}}
\newcommand{\EE}{\mathbb{E}}
\newcommand{\NN}{\mathcal{N}}
\newcommand{\ZZ}{\mathcal{Z}}
\newcommand{\var}{var}
\newcommand{\sd}{sd}

\setlength{\parindent}{0pt}
\renewcommand{\baselinestretch}{1.25}

\pagestyle{fancy}
\fancyhf{}
\renewcommand{\headrule}{\color{structurecolor}\hrule width\textwidth}
\newcommand{\chapterlabel}{}
\fancyhead[L]{\small 生物統計\quad \chapterlabel}
\fancyhead[R]{\small\handoutkind}
\fancyfoot[C]{\small\thepage}

% \grouptitle{組別}{難度星數}{主題}
\newcommand{\grouptitle}[3]{%
    {\Large\bfseries\color{structurecolor} 第 #1 組\quad #3}\hfill{\color{third}#2}\par\medskip}

% Scenario box
\newtcolorbox{scenario}{enhanced, breakable, colback=structurecolor!5, colframe=structurecolor,
    boxrule=0.6pt, arc=2pt, left=6pt, right=6pt, top=4pt, bottom=4pt,
    fonttitle=\bfseries, title=情境}

\newcommand{\stars}[1]{\small\count255=0 \loop\ifnum\count255<#1 $\bigstar$\advance\count255 by 1 \repeat}

% \scenariotitle{編號}{主題}{章節}: heading for a multi-part scenario
\newcommand{\scenariotitle}[3]{%
    {\Large\bfseries\color{structurecolor} 情境 #1\quad #2}\hfill{\color{third}\small #3}\par\medskip}

% Numbered subproblems with a difficulty marker: \sub{星數} text
\newenvironment{subproblems}{\begin{enumerate}[leftmargin=*, label=\textbf{\arabic*.}, itemsep=8pt]}{\end{enumerate}}
\newcommand{\sub}[1]{\item {\color{third}\stars{#1}}\quad}

\renewcommand{\chapterlabel}{第 5--6 章\ 抽樣分布、中央極限定理與區間估計}
\newcommand{\handoutkind}{課堂練習\ 參考解答}

\newenvironment{answers}{\begin{enumerate}[leftmargin=*, label=\textbf{\arabic*.}, itemsep=6pt]}{\end{enumerate}}

\begin{document}

%% ---------------------------------------------------------------- 情境 1
\scenariotitle{1}{社區高血壓篩檢}{第 5、6 章}
\begin{answers}
    \item 是：$\bar{X}$ 由樣本算出且不含未知參數。$\EE(\bar{X}) = \mu$；標準誤 $\sigma/\sqrt{n} = \sigma/\sqrt{64} = \sigma/8$。
    \item 無法\textbf{精確}計算：不知道母體分布，就無法寫出 $\bar{X}$ 確切的抽樣分布。但可以\textbf{近似}計算：中央極限定理只需要樣本獨立同分布、母體平均與變異數存在（題幹已給），$n = 64$ 夠大時 $\bar{X} \leadsto \NN(\mu, \sigma^2/64)$，
          $\PP(\bar{X} > \mu + \sigma/4) = \PP\Big(\frac{\bar{X}-\mu}{\sigma/8} > \frac{\sigma/4}{\sigma/8}\Big) \approx \PP(\ZZ > 2) = 0.0228$。不需要知道 $\mu$、$\sigma$ 的數值。\\
          單一成人則不行：中央極限定理只描述樣本平均的分布；單一觀察值的分布形狀未知，無法計算。
    \item 平均收縮壓：母體常態、$\sigma$ 未知，用 $t$ 區間。標準誤 $15.2/\sqrt{64} = 1.9$；$131.2 \pm 1.998 \times 1.9 = (127.4, 135.0)$ mmHg。$n = 64$ 時改用 $z_{0.025} = 1.96$ 得 $(127.5, 134.9)$，幾乎相同。\\
          盛行率：$\hat{p} = 19/64 = 0.297$，標準誤 $\sqrt{0.297 \times 0.703/64} = 0.057$；\\
          $0.297 \pm 1.96 \times 0.057 = (0.185, 0.409)$。$n\hat{p} = 19$、$n(1-\hat{p}) = 45$ 均大於 5，條件滿足。
    \item 問題在於把\textbf{母體平均}的信賴區間當成\textbf{個別成人}的分布範圍。信賴區間反映的是估計 $\mu$ 的不確定性，寬度由標準誤 $s/\sqrt{n}$ 決定；個別成人之間的差異則由標準差 $\sigma$ 決定。\\
          常態母體約有 $95\%$ 的個體落在 $\mu \pm 1.96\sigma$，以樣本估計值代入：$131.2 \pm 1.96 \times 15.2 = (101.4, 161.0)$ mmHg，遠寬於 $(127.4, 135.0)$。（這是近似，未計入 $\mu$、$\sigma$ 本身的估計誤差。）\\
          樣本數增為 256 人：信賴區間半寬約為 $1.96 \times 15.2/\sqrt{256} \approx 1.86$ mmHg，約減為原來的一半（$\sqrt{256/64} = 2$）；個別成人的範圍則不會隨樣本數變窄，仍約為 $\pm 30$ mmHg，樣本數增加只會讓 $\mu$、$\sigma$ 的估計更準確。
\end{answers}
\newpage

%% ---------------------------------------------------------------- 情境 2
\scenariotitle{2}{急診候診時間改善方案}{第 5、6 章}
\begin{answers}
    \item $\PP(X < 0) = \PP\big(\ZZ < \frac{0-40}{40}\big) = \PP(\ZZ < -1) = 0.1587$。等候時間不可能為負，所以常態分布不適合；平均值等於標準差、又不能小於 0，表示分布明顯右偏（多數人等得不久，少數人等很久）。
    \item 標準誤 $40/\sqrt{100} = 4$；$\PP(\bar{X} < 32) \approx \PP(\ZZ < -2) = 0.0228$。\\
          $n = 100$ 夠大，由中央極限定理，不論母體分布形狀為何，樣本平均都近似常態分布。
    \item 由卜瓦松分布的時間切割性質，30 天的平均總人數為 $4 \times 30 = 120$，故 $T \sim Poisson(120)$。$\lambda > 10$，可近似為 $\NN(120, 120)$。\\
          $\PP(T \ge 140) \approx \PP(T \ge 139.5) = \PP\big(\ZZ \ge \frac{139.5-120}{\sqrt{120}}\big) = \PP(\ZZ \ge 1.78) \approx 0.0375$。（精確的卜瓦松機率約為 0.040。）
    \item 標準誤 $12/\sqrt{16} = 3$。\\
          雙尾：$31 \pm 2.131 \times 3 = (24.6, 37.4)$ 分鐘。\\
          左尾：$(-\infty, 31 + 1.753 \times 3) = (-\infty, 36.3)$ 分鐘。\\
          兩個區間都不包含 40 分鐘，暗示改善後的平均等候時間低於過去（第 7 章會把這個想法正式化為假說檢定）。
\end{answers}

\end{document}
